\section{Density of the actual relative frames and the limiting derivative graph}
\label{sec:nested-relative-frame-density}

This working continuation proves the nested-frame limit for the original
arithmetic amplitude. The input is the sum-connection construction
(FC.1)--(FC.36), retained at the source hash recorded in the accompanying
receipt. All statements below use the original measure, full packet,
relative coordinates, polynomial source, and derivative. The proof uses
an exponential moment that is established directly from the arithmetic
amplitude. No recurrence or Christoffel asymptotic is assumed.

\subsection{An exponential bound with the original arithmetic factors}

Let
\[
 \begin{gathered}
 g(s)=2\xi(s)=s(s-1)\pi^{-s/2}\Gamma(s/2)\zeta(s),\qquad
 h(s)=\prod_{\rho\in\mathcal Z}(s-\rho)^{m_\rho},\\
 d=\deg h,\quad v_h(s)=g(s)/h(s),\qquad
 a(t)=\frac{v_h(1/2+it)}{\sqrt{2\pi}},\quad
 w(t)=|a(t)|^2,\quad \mu_h=\int_{\mathbb R}w(t)\,dt.
 \end{gathered}
 \tag{ND.1}
\]
Here the set of selected zeros is finite and each selected multiplicity
is the full multiplicity in \(g\). The quotient is its entire extension
at those zeros. The empty product \(h=1\) is allowed. The packet is
stable under \(\rho\mapsto1-\overline\rho\), as in FC.1. For a
monic polynomial this says
\(h(1-\overline s)=(-1)^d\overline{h(s)}\); the corresponding
identity for \(g\) has factor one. Hence on the original line
\[
 \overline{a(t)}=(-1)^d a(t),\qquad
 \overline{a^{(j)}(t)}=(-1)^d a^{(j)}(t).
 \tag{ND.2}
\]
There is no change of phase or division by the mass \(\mu_h\).

\begin{theorem}[Actual amplitude and all derivative tails]
For every integer \(j\ge0\) and every real \(b\) with
\(0<b<\pi/4\), there is a finite constant \(C_{h,j,b}\) such that
\[
 |a^{(j)}(t)|\le C_{h,j,b}e^{-b|t|}\qquad(t\in\mathbb R).
 \tag{ND.3}
\]
The amplitude is a nonzero entire function of \(t\). In particular
\(0<\mu_h<\infty\), all polynomially weighted derivative integrals
used below exist, and its real zero set is discrete.
\end{theorem}

\begin{proof}
We give the tail argument, including the gamma factor. For
\(\lambda>0\), \(\tau\in\mathbb R\), and
\(0<\phi<\pi/2\), rotate Euler's gamma integral in the sector with
angle \(\theta=\operatorname{sgn}(\tau)\phi\), taking \(\theta=0\)
when \(\tau=0\). The principal power is holomorphic on the open
sector. On the circle of radius \(R\) the exponential has modulus at
most \(e^{-R\cos\phi}\); its product with the power and arc length
tends to zero as \(R\to\infty\). The small circular contribution
is bounded by a constant times \(\varepsilon^\lambda\) and tends
to zero. Cauchy's theorem therefore gives
\[
 \begin{aligned}
 \Gamma(\lambda+i\tau)
 &=e^{i\theta(\lambda+i\tau)}
       \int_0^\infty r^{\lambda+i\tau-1}e^{-e^{i\theta}r}\,dr,\\
 |\Gamma(\lambda+i\tau)|
 &\le e^{-\phi|\tau|}\Gamma(\lambda)(\cos\phi)^{-\lambda}.
 \end{aligned}
 \tag{ND.4}
\]
At \(\tau=0\) the second bound follows directly from Euler's
integral. Its constant is bounded when \(\lambda\) ranges over a
compact subinterval of \((0,\infty)\).

For completeness, summing the counting function \(\lfloor x\rfloor\)
on \(\Re s>1\) gives
\(\zeta(s)=s\int_1^\infty\lfloor x\rfloor x^{-s-1}\,dx\).
Splitting \(\lfloor x\rfloor=x-\{x\}\) yields
\[
 \zeta(s)=\frac{s}{s-1}
           -s\int_1^\infty\{x\}x^{-s-1}\,dx.
 \tag{ND.5}
\]
The last integral is absolutely and locally uniformly convergent on
\(\Re s>0\); its differentiated integrals converge on compact
subsets because \(\int_1^\infty x^{-1-\sigma}(\log x)^n dx\)
is finite for every \(\sigma>0\). Thus (ND.5) continues the same
zeta function to that half-plane away from \(s=1\).
For \(s=\sigma+it\), \(1/4\le\sigma\le3/4\), we have
\(|s/(s-1)|\le5\), \(|s|/\sigma\le1+4|t|\), and consequently
\(|\zeta(s)|\le6+4|t|\).

Write \(z=t+iv\) with \(|v|\le1/4\), so that
\(s=1/2+iz=(1/2-v)+it\). In (ND.4) take
\(\lambda=(1/2-v)/2\in[1/8,3/8]\) and \(\tau=t/2\).
The modulus of \(\pi^{-s/2}\) is exactly \(\pi^{-(1/2-v)/2}\).
The two polynomial factors in \(g\), followed by (ND.5), therefore
give the bound
\[
 |g(1/2+iz)|\le C_\phi(1+|t|)^3 e^{-\phi|t|/2}
       \qquad(|\Im z|\le1/4).
 \tag{ND.6}
\]
For \(|t|\ge\max(1,2\max_{\rho\in\mathcal Z}|\Im\rho|)\),
each factor \(|s-\rho|\) is at least \(|t|/2\). Thus
\(|h(s)|\ge2^{-d}|t|^d\); for the empty product the lower bound
is one and the maximum over roots is unnecessary. Division by this
original polynomial, without removing its remaining unit, bounds
\(v_h\) on the unbounded part of the strip. The complementary
closed rectangle is compact and the entire quotient is bounded there,
including every cancelled root. Choose \(2b<\phi<\pi/2\).
The polynomial factor times \(e^{-(\phi/2-b)|t|}\) is bounded.
It follows that
\(|a(z)|\le C_{h,b}e^{-b|\Re z|}\) on this strip, with the
original factor \(1/\sqrt{2\pi}\) included in \(C_{h,b}\).
The Cauchy estimate on the circle of radius \(1/8\) about a real
\(t\) now gives
\[
 |a^{(j)}(t)|\le j!8^j C_{h,b}e^{b/8}e^{-b|t|}.
\]
This proves (ND.3). Entirety follows from cancellation at the specified
finite roots. The function cannot vanish identically: the formula for
\(g\) is nonzero on the real interval \(s>1\), and division by a
nonzero polynomial preserves nonidentity. The identity theorem makes
the real zeros isolated. Positivity and finiteness of the mass follow
from nonidentity and (ND.3).
\end{proof}

\subsection{Every centred fibre and its exponential moment}

Fix \(k\ge2\), put \(r=k-1\), and retain the real coordinates
\[
 \begin{gathered}
 u=\sum_{i=1}^k t_i,\quad y_i=t_i-u/k\ (1\le i\le r),\quad
 y_k=-\sum_{i=1}^r y_i,\quad t_i=u/k+y_i,\\
 S=k/2+iu,\quad z_i=s_i-S/k=iy_i,\qquad
 \Psi_u(y)=\prod_{i=1}^k a(u/k+y_i),\quad
 d\mu_u(y)=|\Psi_u(y)|^2\,dy.
 \end{gathered}
 \tag{ND.7}
\]
The independent relative variables are \(y_1,\ldots,y_r\).
The determinant of the map with ordered variables
\((u,y_1,\ldots,y_r)\mapsto(t_1,\ldots,t_k)\) is
\((-1)^r\). Indeed moving the first column to the last contributes
\((-1)^r\), and the resulting block determinant is
\(\det\left(\begin{smallmatrix}I_r&\mathbf1/k\\
-\mathbf1^T&1/k\end{smallmatrix}\right)=1/k+r/k=1\).
In particular the original product measure is exactly \(du\,dy\).
For \(k=2\), \(v=t_1-t_2=2y_1\) gives \(du\,dv/2\).
Neither the centre \(k/2\) nor this Jacobian is changed.

For every fixed \(u\), each of the \(k\) affine functions
\(y\mapsto u/k+y_i\) is nonconstant. The inverse image of a
real zero of \(a\) is a proper affine hyperplane. There are at
most countably many such hyperplanes, each of Lebesgue measure zero.
Consequently \(\Psi_u\ne0\) almost everywhere for every \(u\),
including any fibre on which a particular normal derivative vanishes.
The orders of the zeros remain the orders of the original product;
the measure assertion does not modify the product on a hyperplane.

For every \(0<\eta<\pi/2\), choose \(b\) so that
\(\eta<2b<\pi/2\). Since
\(\sum_{i=1}^k|t_i|\ge\|y\|_1-r|u|/k\), (ND.3) proves
the explicit bound
\[
 \begin{aligned}
 \int_{\mathbb R^r}e^{\eta\|y\|_1}\,d\mu_u(y)
 &\le C_{h,0,b}^{2k}e^{2br|u|/k}
       \int_{\mathbb R^r}e^{-(2b-\eta)\|y\|_1}\,dy\\
 &=C_{h,0,b}^{2k}e^{2br|u|/k}
       \left(\frac{2}{2b-\eta}\right)^r<\infty.
 \end{aligned}
 \tag{ND.8}
\]
This bound is uniform on compact \(u\)-intervals. It proves both
finite moments and \(0<m_k(u):=\mu_u(\mathbb R^r)<\infty\)
for every \(u\). The coordinate change and Tonelli's theorem retain
the full mass identity
\[
 \int_{\mathbb R}m_k(u)\,du=\mu_h^k.
 \tag{ND.9}
\]
At fixed relative coordinate the exact derivative is
\[
 \Psi'_u(y)=\frac1k\sum_{i=1}^k a'(u/k+y_i)
                         \prod_{j\ne i}a(u/k+y_j).
 \tag{ND.10}
\]
Every derivative of this product, multiplied by any polynomial in
\((u,y)\), is square integrable in \(du\,dy\). To verify this
claim directly, every derivative is a finite sum of products of
derivatives of the original \(a(t_i)\), with the powers of \(1/k\)
from (ND.7) retained. Each product is bounded by a constant times
\(e^{-b\sum|t_i|}\) by (ND.3). A polynomial in \((u,y)\)
becomes a polynomial in \(t_1,\ldots,t_k\) under the displayed
invertible linear map; each of its monomials is integrable against
\(e^{-2b\sum|t_i|}\). The same bound on a compact \(u\)-interval
proves continuity, and all orders of differentiation, of the relevant
functions with values in \(L^2(dy)\), by dominated convergence.

\subsection{A complete weighted polynomial density proof}

\begin{theorem}[Exponential moments imply the required density]
Let \(\mu\) be a finite positive Borel measure on \(\mathbb R^r\)
such that \(\int e^{\eta\|y\|_1}d\mu(y)<\infty\) for some
\(\eta>0\). Complex polynomials in all \(r\) coordinates are
dense in \(L^2(\mu)\).
\end{theorem}

\begin{proof}
All polynomials belong to this space by the exponential moment. Suppose
\(f\in L^2(\mu)\) is orthogonal to every polynomial, using inner
products conjugate linear in the first entry, and define the finite
complex measure \(d\nu=\overline f\,d\mu\). For \(0\le2\tau<\eta\),
Cauchy--Schwarz gives
\[
 \int e^{\tau\|y\|_1}d|\nu|(y)
 \le\|f\|_{L^2(\mu)}
       \left(\int e^{2\tau\|y\|_1}d\mu(y)\right)^{1/2}<\infty.
 \tag{ND.11}
\]
Every monomial moment of \(\nu\) is zero. Fix
\(\xi\in\mathbb R^r\setminus\{0\}\). The function
\[
 F_\xi(z)=\int e^{iz\xi\cdot y}\,d\nu(y)
\]
is holomorphic on the connected strip
\(|\Im z|\|\xi\|_\infty<\eta/2\). To see the differentiability,
restrict to a compact substrip and choose \(\tau<\eta/2\) strictly
larger than its bound on \(|\Im z|\|\xi\|_\infty\). Every
polynomial factor from differentiation is bounded by the remaining
exponential slack in (ND.11), which justifies differentiation under
the integral. In a disc about zero, the exponential series is
absolutely integrable by the same bound. Its \(n\)-th coefficient
is \(i^n\int(\xi\cdot y)^n d\nu(y)/n!=0\), since this is a
polynomial moment. Thus \(F_\xi\) vanishes near zero and throughout
the strip by the holomorphic identity theorem. In particular
\(\widehat\nu(\xi)=F_\xi(1)=0\). The zero moment also gives
\(\widehat\nu(0)=0\).

Here is the Fourier uniqueness argument needed for this finite complex
measure. For \(\varepsilon>0\) put
\[
 G_\varepsilon(x)=(4\pi\varepsilon)^{-r/2}
                  e^{-|x|^2/(4\varepsilon)}
  =(2\pi)^{-r}\int_{\mathbb R^r}
                  e^{-\varepsilon|\xi|^2}e^{i\xi\cdot x}\,d\xi.
 \tag{ND.12}
\]
The equality follows in one variable by differentiating the Gaussian
Fourier integral and integrating by parts: its derivative is
\(-x/(2\varepsilon)\) times itself and its value at zero is
\(\sqrt{\pi/\varepsilon}\). The latter Gaussian integral follows
by squaring it and using polar coordinates. Taking products gives
(ND.12). Since \(\nu\) is finite and the frequency Gaussian is
integrable, Fubini's theorem gives, with the same positive Fourier sign,
\[
 (G_\varepsilon*\nu)(x)=(2\pi)^{-r}\int_{\mathbb R^r}
 e^{-\varepsilon|\xi|^2}e^{i\xi\cdot x}\widehat\nu(-\xi)\,d\xi=0.
\]
For every bounded uniformly
continuous function \(b\), another absolutely convergent Fubini
integral yields \(\int(b*G_\varepsilon)d\nu=0\). Moreover
\(b*G_\varepsilon\to b\) uniformly: in its difference from \(b\),
the integration over \(|x|<\delta\) is bounded by the uniform
modulus of continuity at \(\delta\), while the complementary
integral is at most \(2\|b\|_\infty\int_{|x|\ge\delta}G_\varepsilon\),
which tends to zero after the substitution \(x=\sqrt\varepsilon z\).
Thus \(\int b\,d\nu=0\) for all such \(b\).

For each nonempty closed set \(F\subset\mathbb R^r\), the bounded
Lipschitz functions
\(b_n(x)=\max(0,1-n\operatorname{dist}(x,F))\) converge pointwise
to \(1_F\). Dominated convergence against \(|\nu|\) proves
\(\nu(F)=0\); this also holds for the empty closed set. In particular
\(\nu(\mathbb R^r)=0\). The Borel sets with zero \(\nu\)-measure
form a Dynkin system: complements have measure zero by this total
mass identity, and countable disjoint unions have measure zero by
countable additivity. Closed sets form a pi-system generating the
Borel sigma-algebra. The pi--lambda theorem therefore gives
\(\nu(A)=0\) for every Borel set \(A\). Applying this to the
positive and negative level sets of \(\Re f\) and \(\Im f\)
shows \(f=0\) almost everywhere. Finally, the orthogonal complement
of the closed polynomial span is zero, so the Hilbert orthogonal
decomposition proves that span is all of \(L^2(\mu)\).
\end{proof}

\subsection{The multiplier, literal frames, and their strong limit}

Let \(\mathscr H=L^2(\mathbb R^r,dy)\). For each fixed \(u\),
the exact multiplication map is
\[
 \begin{aligned}
 \mathcal M_u:L^2(\mu_u)&\longrightarrow\mathscr H,
       &q&\longmapsto\Psi_u q,\\
 \mathcal M_u^{-1}F(y)&=
 \begin{cases}F(y)/\Psi_u(y),&\Psi_u(y)\ne0,\\
 0,&\Psi_u(y)=0.
 \end{cases}
 \end{aligned}
 \tag{ND.13}
\]
It is a unitary map: its squared norm is exactly
\(\int|q|^2|\Psi_u|^2dy\), and the inverse has weighted squared
norm \(\int|F|^2dy\) because the zero set is null. These are maps
between the displayed weighted and unweighted spaces. No claim of
bounded multiplication by \(1/\Psi_u\) on unweighted
\(\mathscr H\) is made or used. Assigning zero on null hyperplanes
only selects a representative of a Hilbert class; every analytic
amplitude and all of its pointwise zero orders remain those of (ND.7).

For \(m\ge0\) use the full monomial family
\[
 \begin{gathered}
 \mathcal A_m=\{\alpha\in\mathbb N^r:|\alpha|\le m\},\quad
 \theta_\alpha(z)=z^\alpha,\quad
 \theta_\alpha(iy)=i^{|\alpha|}y^\alpha,\\
 j_m(u):\mathbb C^{\mathcal A_m}\to\mathscr H,\quad
 j_m(u)c=\Psi_u\sum_{\alpha\in\mathcal A_m}
                         i^{|\alpha|}y^\alpha c_\alpha,\\
 V_m(u)=\operatorname{ran}j_m(u),\quad W_m(u)=j_m(u)^*j_m(u),
 \quad \Pi_m(u)=j_m(u)W_m(u)^{-1}j_m(u)^*.
 \end{gathered}
 \tag{ND.14}
\]
All moments defining these matrices are finite by (ND.8). If \(c\ne0\),
its displayed polynomial is nonzero. By continuity it is nonzero in a
nonempty open set, which has positive \(\mu_u\)-measure because
the density is positive almost everywhere. Thus \(W_m(u)>0\)
for every \(u\). Multiplication proves that \(\Pi_m\) is
self-adjoint, idempotent, and has range \(V_m\), so it is precisely
the orthogonal projection onto that range. The inverse here is the
positive source Gram \(W_m^{-1}\); no inverse of a possibly zero
normal Gram is used.

For \(n\ge m\), let \(E_{nm}\) be literal zero-padding of the
coefficient array. The complete inclusion and Gram maps are
\[
 \begin{gathered}
 j_n E_{nm}=j_m,\quad W_m=E_{nm}^*W_n E_{nm},\quad
 E_{\ell n}E_{nm}=E_{\ell m},\\
 V_m\subset V_n,\qquad
 \Pi_n\Pi_m=\Pi_m\Pi_n=\Pi_m.
 \end{gathered}
 \tag{ND.15}
\]
The first line follows term by term in the original polynomial sum.
The second line follows from the nested ranges and orthogonal
complements. In particular \(\Pi_n-\Pi_m\) is the orthogonal
projection onto \(V_n\cap V_m^\perp\).

\begin{theorem}[Full relative frames on every actual fibre]
For every \(u\in\mathbb R\) and every \(F\in\mathscr H\),
\[
 \lim_{m\to\infty}\|F-\Pi_m(u)F\|_{\mathscr H}=0.
 \tag{ND.16}
\]
For every finite \(m\), however,
\(\|I-\Pi_m(u)\|_{\mathscr H\to\mathscr H}=1\).
\end{theorem}

\begin{proof}
The density theorem applies to \(\mu_u\) by (ND.8), separately
for each real \(u\). The phases \(i^{|\alpha|}\) are nonzero
constants and hence the literal monomials in (ND.14) have precisely
the full polynomial span. Applying the onto isometry (ND.13) proves
\(\overline{\bigcup_m V_m(u)}=\mathscr H\). Given \(F\) and
\(\varepsilon>0\), choose \(v\in V_M(u)\) with
\(\|F-v\|<\varepsilon\). For \(m\ge M\), orthogonal
projection minimizes distance, so
\(\|F-\Pi_mF\|\le\|F-v\|<\varepsilon\).
This proves (ND.16).
The Hilbert space \(L^2(\mathbb R^r)\) is infinite dimensional:
indicators of infinitely many disjoint finite positive-measure balls
are linearly independent. Each \(V_m\) has finite dimension.
Its orthogonal complement therefore contains a unit vector, on which
\(I-\Pi_m\) has norm one. Orthogonal projection contraction gives
the opposite inequality, proving the asserted operator norm exactly.
\end{proof}

\subsection{A fixed original source and the complete derivative domain}

Fix one original polynomial \(P(s_1,\ldots,s_k)\) of total degree
\(D\). In the exact linear coordinates
\(S=\sum s_i\), \(z_i=s_i-S/k\), \(\sum z_i=0\), it has
the unique expansion
\[
 P=\sum_{|\alpha|\le D}z^\alpha f_\alpha(S),\qquad
 \deg f_\alpha\le D-|\alpha|.
 \tag{ND.17}
\]
Uniqueness is coefficient comparison in the independent variables
\((S,z_1,\ldots,z_r)\); substitution by an invertible linear
map preserves total degree. For \(m\ge D\), retain an admitted
source degree \(M_m\ge\max(D,m)\) and the full relative family
\(\mathcal A_m\). The fixed source has coefficient column
\[
 c_{P,m}(u)_\alpha=
 \begin{cases}f_\alpha(k/2+iu),&|\alpha|\le D,\\
 0,&D<|\alpha|\le m.
 \end{cases}
 \qquad c_{P,n}=E_{nm}c_{P,m}.
 \tag{ND.18}
\]
Every original constraint \(\deg f_\alpha\le M_m-|\alpha|\)
is respected; neither the fixed polynomial nor its coefficient degree
is changed as the frame grows. One may in particular take \(M_m=m\)
for all \(m\ge D\), which is the full degree-\(m\) source of FC.28.

Put
\[
 \begin{gathered}
 p_u(y)=P(k/2+iu,iy_1,\ldots,iy_r),\quad
 A_P(u,y)=\Psi_u(y)p_u(y)=j_m(u)c_{P,m}(u),\\
 \partial_S^{\rm rel}=\frac1k\sum_{i=1}^k\partial_{s_i},\quad
 F_P(u,y)=\Psi'_u(y)p_u(y),\\
 \partial_u A_P=F_P+i\Psi_u(\partial_S^{\rm rel}P)_u.
 \end{gathered}
 \tag{ND.19}
\]
The notation \(P(S,z)\) in these formulas denotes the exact
substitution \(s_i=S/k+z_i\); \(z_k=-\sum_{i<k}z_i\).
At fixed relative coordinates \(dS/du=i\), so the displayed
derivative includes exactly the factor \(i/k\) on the sum of
the original coordinate derivatives. By (ND.3) and (ND.10),
\(A_P,F_P,\partial_uA_P\) belong to
\(\mathscr K=L^2(\mathbb R_u;\mathscr H)\).
In fact all their polynomially weighted derivatives belong to
\(\mathscr K\). In particular \(A_P\) lies in the actual
derivative domain
\[
 \mathcal D_u=\{A\in\mathscr K:\partial_uA\in\mathscr K
                         \text{ in the distributional sense}\}.
 \tag{ND.20}
\]
The smooth derivative in (ND.19) equals the distributional derivative,
as follows by integration by parts against compactly supported test
functions and the preceding integrability bounds.

We also retain its original Mellin realization and sign. If
\(\mathcal M f(s)=\int_0^\infty f(x)x^{s-1}dx\), the original
test vector \(F_h\) has \(\mathcal M F_h=v_h\), and the
polynomial source is
\(\mathcal T_h^{(k)}P=P(D_1,\ldots,D_k)F_h^{\otimes k}\),
where \(D_i=-x_i\partial_{x_i}\). The normalized critical-line
transform has factor \((2\pi)^{-k/2}\), followed by (ND.7).
Integration by parts gives \(\mathcal M(D_i f)=s_i\mathcal M f\),
so its image is exactly \(A_P\), with the original product of
\(a\)'s. Writing \(x_i=e^{q_i}\), the Mellin integral becomes
the Fourier transform with kernel \(e^{it_iq_i}\) of
\(e^{(q_1+\cdots+q_k)/2}f(e^{q_1},\ldots,e^{q_k})\).
The change of variable has squared norm exactly \(\int|f|^2d\mathbf x\);
Fourier Plancherel and the Jacobian one in (ND.7) give the isometry
\(\mathscr U_k\) in FC.9. On this original test domain,
\[
 \begin{gathered}
 \mathscr U_k\mathcal T_h^{(k)}P=A_P,\qquad
 \mathscr U_kD^{(k)}=S\mathscr U_k,\qquad
 \partial_u\mathscr U_k=i\mathscr U_k\mathscr L_k,\\
 D^{(k)}=\sum_iD_i,\qquad
 \mathscr L_k=\frac1k\sum_i\log x_i,\qquad
 [D^{(k)},\mathscr L_k]=-I.
 \end{gathered}
 \tag{ND.21}
\]
The second derivative identity follows by differentiating the Fourier
integral with \(\partial t_i/\partial u=1/k\); the commutator
follows from \(-x_i\partial_{x_i}\log x_i=-1\), summed with
the displayed coefficient \(1/k\). All these operations are on
the original test vectors and their logarithmic derivatives. Their
spectral images have the established \(L^2\) derivatives, so no
new domain or arbitrary assignment of an arithmetic jet is introduced.

In the degree-\(m\) frame define the original matrices and maps
\[
 \begin{gathered}
 B_m=j_m^*j_m',\quad \Gamma_m=W_m^{-1}B_m,\quad
 N_m=(I-\Pi_m)j_m',\quad \mathcal N_m=N_m^*N_m,\\
 t_m(P)=j_m(c_{P,m}'+\Gamma_m c_{P,m}),\qquad
 r_m(P)=N_m c_{P,m}.
 \end{gathered}
 \tag{ND.22}
\]
Every term is the FC.6 term in the literal enlarged frame. Since
\(j_m\Gamma_m=\Pi_mj_m'\), the product rule gives
\[
 \begin{gathered}
 t_m(P)=\Pi_m\partial_uA_P,\qquad
 r_m(P)=(I-\Pi_m)\partial_uA_P=(I-\Pi_m)F_P,\\
 \partial_uA_P=t_m(P)+r_m(P),\quad
 \langle t_m(P,u),r_m(P,u)\rangle_{\mathscr H}=0,\\
 \|\partial_uA_P\|_{\mathscr K}^2
     =\|t_m(P)\|_{\mathscr K}^2+\|r_m(P)\|_{\mathscr K}^2.
 \end{gathered}
 \tag{ND.23}
\]
For the second equality, the coefficient derivative in (ND.19) belongs
to \(V_m(u)\), so its orthogonal complement is zero. This equality
identifies a single fixed function \(F_P\), independent of frame
degree, which controls every residual.

\begin{theorem}[Vanishing of the fixed-source actual normal residual]
For every fixed original polynomial \(P\) and every real \(u\),
\[
 \begin{gathered}
 \|r_m(P,u)\|_{\mathscr H}\longrightarrow0,\qquad
 \|t_m(P,u)-\partial_uA_P(u)\|_{\mathscr H}\longrightarrow0,\\
 \|r_m(P)\|_{\mathscr K}\longrightarrow0,\qquad
 \|t_m(P)-\partial_uA_P\|_{\mathscr K}\longrightarrow0.
 \end{gathered}
 \tag{ND.24}
\]
The fibre convergence of both norms is uniform on compact
\(u\)-intervals. The derivative graph converges, in the original
\(\mathscr K\oplus\mathscr K\) norm, to
\((\partial_uA_P,0)\), while its sum is exactly \(\partial_uA_P\)
at every degree.
\end{theorem}

\begin{proof}
For each \(u\), \(F_P(u)\in\mathscr H\), and (ND.16)
applied to this vector gives the first limit through (ND.23).
Furthermore,
\[
 0\le\|r_m(P,u)\|^2\le\|F_P(u)\|^2,
 \qquad \int_{\mathbb R}\|F_P(u)\|^2du<\infty.
 \tag{ND.25}
\]
All functions here are measurable; in fact \(j_m,j_m',W_m\) are
smooth in \(u\) by the compact-interval bounds after (ND.10),
and \(W_m^{-1}\) is smooth since its determinant is everywhere
positive. The dominated convergence theorem proves the integrated
residual limit. Since \(t_m-\partial_uA_P=-r_m\), it proves
the two tangential limits and the asserted graph convergence as well.

For \(n\ge m\), (ND.15) gives the exact orthogonal decomposition
\[
 \begin{gathered}
 r_m(P,u)=r_n(P,u)+(\Pi_n-\Pi_m)F_P(u),\\
 \|r_m(P,u)\|^2-\|r_n(P,u)\|^2
              =\|(\Pi_n-\Pi_m)F_P(u)\|^2\ge0.
 \end{gathered}
 \tag{ND.26}
\]
The two terms are orthogonal because the first lies in \(V_n^\perp\)
and the second in \(V_n\). Each squared residual is continuous
by the same smoothness argument. On a compact set \(K\) the open
sets \(\{u:\|r_m(P,u)\|^2<\varepsilon\}\) are increasing and
cover \(K\). A finite subcover, followed by taking its largest
index, makes the squared residual less than \(\varepsilon\)
everywhere on \(K\). This proves compact uniform convergence
without a global uniform-in-source assertion.
\end{proof}

For a fixed finite source core \(\operatorname{span}\{P_1,\ldots,P_q\}\)
let \(R_m c=\sum_j c_jr_m(P_j)\). Cauchy--Schwarz gives
\[
 \|R_m\|^2\le\sum_{j=1}^q\|r_m(P_j)\|_{\mathscr K}^2\to0,
 \quad \|R_m^*R_m\|\to0,\quad
 R_m^*R_m-R_n^*R_n\succeq0\ (n\ge m).
 \tag{ND.27}
\]
The last inequality is (ND.26) integrated for each linear combination
of the fixed sources. These are exactly the normal Gram matrices
restricted by literal inclusion of that source core. The core and its
coefficient norm are fixed in this statement. The full degree-varying
source dimension, its minimum arithmetic norm, and the size of its
normal-image inverse are not replaced by this finite-core limit.

The exact translation to the original source-amplitude metric is also
available on this same fixed core. For an independent core put
\((G_0)_{ij}=\langle A_{P_i},A_{P_j}\rangle_{\mathscr K}\).
It is independent of \(m\) and strictly positive: a nonzero linear
combination of the original polynomials is nonzero on an open subset
of the real critical-line coordinates, and the original product
amplitude is nonzero almost everywhere there. Its squared integral
is therefore positive. With \(H_m=R_m^*R_m\), define
\[
 K_m=G_0^{-1}H_m,\quad K_m^*G_0=G_0K_m=H_m,\quad
 \|R_m\|_{G_0\to\mathscr K}^2
 \le\lambda_{\min}(G_0)^{-1}
                \sum_j\|r_m(P_j)\|_{\mathscr K}^2\longrightarrow0.
\]
The first equalities follow by matrix multiplication, and the bound
follows from \(c^*G_0c\ge\lambda_{\min}(G_0)\|c\|^2\) and
the same column Cauchy--Schwarz inequality. Multiplication by the
positive square root \(G_0^{1/2}\) is the isometry from this original
coefficient metric to the Euclidean metric. It carries \(K_m\) to
the positive matrix \(G_0^{-1/2}H_mG_0^{-1/2}\); its largest
eigenvalue is exactly \(\|R_m\|_{G_0\to\mathscr K}^2\).
Thus \(K_m\) also tends to zero in its original metric operator
norm. These explicit congruence and comparison maps preserve both
the source-amplitude metric and the normal metric.

The arithmetic maps themselves are retained by the same inclusions.
For a nonempty packet set
\(I=(h(s_1),\ldots,h(s_k))\), let \(j_I\) denote the full finite
jet map, and put \(U(\mathbf s)=\prod_i v_h(s_i)\).
At every selected root \(v_h(\rho)\ne0\), since its full zero
order has been divided out; hence \(\upsilon=j_IU\) is a unit
in the full tensor jet algebra. The original source and logarithmic
derivative observations are
\[
 J(\mathcal T_h^{(k)}P)=\upsilon[P],\qquad
 J\mathscr L_k\mathcal T_h^{(k)}P
              =j_I\bigl(\partial_S^{\rm rel}(UP)\bigr).
 \tag{ND.28}
\]
Indeed the Mellin transform of \(\mathcal T_h^{(k)}P\) is \(UP\);
insertion of \(k^{-1}\sum\log x_i\) differentiates it by
\(\partial_S^{\rm rel}\). At a selected tuple the algebra of
germs modulo the full powers is finite; a germ with nonzero constant
term is invertible by its finite geometric expansion in the nilpotent
maximal ideal. This proves the asserted unit property including all
mixed jet orders. Since (ND.18) is the identity on \(P\), both
quantities in (ND.28) are identical at every degree. No passage of
this jet map through a Hilbert-space limit is required for, or asserted
by, the density theorem. If \(h=1\), the quotient by \(I=(1)\)
is the zero algebra and both finite arithmetic targets are zero.

\subsection{Invariant frames, sparse frames, and all original zero cases}

Let \(G\) be any fixed subgroup of the factor permutation group
\(S_k\). Its action on \(\sum y_i=0\), expressed in the first
\(r\) coordinates, is real linear with determinant of absolute value
one. To check the determinant, the action on the full \(t\) variables
is a permutation matrix with determinant \(\pm1\); conjugating by
(ND.7) leaves \(u\) fixed and acts linearly on \(y\), so the
remaining block has determinant \(\pm1\). The product \(\Psi_u\)
is unchanged by this action. Therefore
\[
 (\mathsf U_gF)(y)=F(g^{-1}y),\qquad
 \mathsf R_G=\frac1{|G|}\sum_{g\in G}\mathsf U_g
 \tag{ND.29}
\]
defines unitary group actions on both \(\mathscr H\) and
\(L^2(\mu_u)\); they are intertwined by \(\mathcal M_u\).
The adjoint of \(\mathsf U_g\) is \(\mathsf U_{g^{-1}}\), so
\(\mathsf R_G^*=\mathsf R_G\). Counting the \(|G|\) products
\(gh=\ell\) for each \(\ell\) gives \(\mathsf R_G^2=\mathsf R_G\).
Its image is the invariant subspace: its image is fixed under every
group element, and its value on a fixed vector is that vector.
Thus it is exactly the invariant orthogonal projection.

Let \(V_m^G(u)\) be \(\Psi_u\) times all \(G\)-invariant
relative polynomials of degree at most \(m\), using any full
independent basis of that space with its original coefficients.
The degree is preserved by the linear action and by its finite average.
Applying \(\mathsf R_G\) to polynomial approximants in the density
theorem proves density in the invariant weighted subspace. Intertwining
by (ND.13) therefore gives
\[
 \overline{\bigcup_m V_m^G(u)}=\mathscr H^G,\qquad
 \Pi_m^G(u)\xrightarrow[\text{strongly}]{}\mathsf R_G.
 \tag{ND.30}
\]
To verify the second formula for an arbitrary \(F\), decompose it
as \(\mathsf R_GF+(I-\mathsf R_G)F\). The second part is
orthogonal to every \(V_m^G\), and the first is approximated by
their dense union by the same best-approximation argument as (ND.16).
If the original polynomial \(P\) is \(G\)-invariant, then
\(A_P,F_P,\partial_uA_P\) are invariant, since the group fixes
\(u\) and \(\Psi_u\). Its expansion in a full homogeneous
invariant basis retains the bound \(\deg f_\alpha\le D-\ell_\alpha\):
apply coefficient comparison in powers of \(S\) and the homogeneous
relative grading to (ND.17). Consequently (ND.23)--(ND.27) hold
for these actual invariant frames, with \(\Pi_m^G\) in place of
\(\Pi_m\). In particular their fixed-source residual also tends
to zero pointwise, on compact \(u\)-intervals, and in \(\mathscr K\).
This includes full symmetric orbit frames without changing the
unscaled orbit coefficients.

For an arbitrary nested family of finite polynomial spaces
\(E_m\), let \(V_m^E(u)=\mathcal M_uE_m\) and
\(V_\infty^E(u)=\overline{\bigcup_m V_m^E(u)}\). Its exact
strong projection limit is
\[
 \Pi_m^E(u)\xrightarrow[\text{strongly}]{}
           Q_u^E:=\operatorname{proj}_{V_\infty^E(u)}.
 \tag{ND.31}
\]
Here is a proof requiring no density assumption. For \(n\ge m\),
orthogonality of the nested projections gives
\(\|\Pi_n^EF-\Pi_m^EF\|^2=\|\Pi_n^EF\|^2-\|\Pi_m^EF\|^2\).
The squared norms increase and are bounded by \(\|F\|^2\).
Thus the projections applied to \(F\) are Cauchy. Their limit
belongs to \(V_\infty^E\), and its difference from \(F\) is
orthogonal to each \(V_m^E\), hence to their closure. This proves
(ND.31) and identifies the exact projection map. A fixed source
admitted from some level onward is a sum of fixed members of that
relative polynomial space times polynomials in \(S\). Its coefficient
derivative replaces each such polynomial by \(i\) times its
\(S\)-derivative and therefore belongs to the same relative space.
The residual limit is consequently precisely
\((I-Q_u^E)\partial_uA_P=(I-Q_u^E)F_P\). This convergence is
also in \(\mathscr K\), by pointwise strong convergence and the
integrable bound \(4\|F_P(u)\|^2\) for the squared difference
from its limiting residual. The limit field is measurable as a
pointwise norm limit of measurable projection fields.

The projection in (ND.31) need not be the identity. A concrete actual
example uses \(E_m=\mathbb C\) at every level. Set
\(b_u=m_k(u)^{-1}\int y_1d\mu_u(y)\), which is finite by (ND.8).
The nonzero vector \(F=\Psi_u(y)(y_1-b_u)\) has zero inner product
with \(\Psi_u\), so every \(\Pi_m^EF=0\). It is nonzero because
\(\mu_u\) has positive density almost everywhere and the affine
hyperplane \(y_1=b_u\) is null. This gives the exact projection
and an explicit failed approximation, without asserting that the
sparse space is unrelated to the full one.

All zero and empty cases have concrete maps. The zero source has
\(A_P=F_P=t_m=r_m=0\) at every degree. A nonzero source whose
arithmetic class \([P]\) or \(\upsilon[P]\) vanishes is still the
same nonzero polynomial in (ND.17), with the same amplitude and
convergence; (ND.28) retains its full arithmetic class. The packet
\(h=1\) gives the actual \(2\xi/\sqrt{2\pi}\) amplitude and
satisfies every analytic estimate above. No actual fibre has zero
mass for \(k\ge2\), by the hyperplane argument following (ND.7).
One may append an initial empty frame \(V_{-1}=\{0\}\) with
\(\Pi_{-1}=0\); the full nonempty sequence starts with the constant
frame at degree zero. Empty source arrays have the unique zero maps.

The original exceptional two-factor fibre is retained. If the packet
makes \(a\) even, as for the quartet in FC.27, then \(a'\) is odd
and (ND.10) gives
\[
 \Psi'_0(y)=\tfrac12\bigl(a'(y)a(-y)+a(y)a'(-y)\bigr)=0,
 \quad j_m'(0)=0,\quad N_m(0)=0,\quad r_m(P,0)=0.
 \tag{ND.32}
\]
The derivative of the coefficient polynomial may be nonzero; (ND.23)
keeps that term in the tangential component. The density proof and
the limiting derivative graph require no positive inverse for
\(\mathcal N_m(0)\). Their only fibre inverse is \(W_m^{-1}\),
which remains strictly positive there. No conclusion about one-factor
relative frames is hidden in the assumption \(k\ge2\): for \(k=1\)
the relative coordinate space has one point, the multiplication image
is one-dimensional when \(a(u)\ne0\) and is zero when \(a(u)=0\).
At such a zero it is not onto the scalar ambient Hilbert space, and
the pointwise density-to-identity assertion (ND.16) was explicitly
proved only for \(k\ge2\).

\subsection{Frame covariance and the exact half-line transport}

Let \(C_m(u)\) be any smooth invertible change of relative frame.
Use \(\widetilde j_m=j_m C_m\) and
\(\widetilde c_{P,m}=C_m^{-1}c_{P,m}\). The product and projection
remain the same, while the complete connection and normal formulas are
\[
 \begin{gathered}
 \widetilde W_m=C_m^*W_mC_m,\quad
 \widetilde B_m=C_m^*B_mC_m+C_m^*W_mC_m',\\
 \widetilde\Gamma_m=C_m^{-1}\Gamma_mC_m+C_m^{-1}C_m',\quad
 \widetilde N_m=N_mC_m,\quad
 \widetilde N_m\widetilde c_{P,m}=r_m(P),\\
 \widetilde j_m(\widetilde c_{P,m}'
                +\widetilde\Gamma_m\widetilde c_{P,m})=t_m(P),
 \quad \widetilde E_{nm}=C_n^{-1}E_{nm}C_m.
 \end{gathered}
 \tag{ND.33}
\]
The first two lines follow by differentiating \(j_mC_m\) and
using the unchanged projection image; its complementary projection
annihilates \(j_mC_m'\). Differentiating
\(C_m^{-1}c_{P,m}\) cancels the \(C_m^{-1}C_m'\) term and
proves the tangential identity. Finally
\(\widetilde j_n\widetilde E_{nm}=\widetilde j_m\), so these
are the exact typed inclusion maps for the same nested images.
Arbitrary growth of \(C_m\) is not used to assign an independent
coefficient domain: the admitted columns are precisely the pullbacks
of (ND.18), and the actual amplitudes and graph norms remain unchanged.

For clarity the residual limit can also be transported into the original
two-branch half-line representation without an evenness assumption.
For \(R\in\mathscr K\), let \(x=u^2>0\), \(v=\sqrt x\), and set
\[
 \begin{gathered}
 R_+(x)=R(v),\quad R_-(x)=R(-v),\quad
 e_R(x)=\tfrac12(R_+(x)+R_-(x)),\quad
 o_R(x)=\frac{R_+(x)-R_-(x)}{2i\sqrt x},\\
 R_+=e_R+i\sqrt x\,o_R,\qquad R_-=e_R-i\sqrt x\,o_R,\\
 \|R\|_{\mathscr K}^2
 =\int_0^\infty\frac{\|R_+\|^2+\|R_-\|^2}{2\sqrt x}\,dx
 =\int_0^\infty\left(\frac{\|e_R\|^2}{\sqrt x}
                        +\sqrt x\,\|o_R\|^2\right)dx.
 \end{gathered}
 \tag{ND.34}
\]
The first equality is the positive Jacobian on each branch; the second
is the parallelogram identity after the displayed exact inverse.
These formulas define an onto isometry between the full-line space
and the displayed two weighted half-line spaces. Thus (ND.24) carries
the actual residual to zero in the sum of both original weighted
norms, and carries the full derivative graph to the same transported
derivative. Pointwise fibre convergence holds at every \(x>0\) on
both branches. At the meeting point \(u=0\), the original smooth
amplitudes have their original trace; no new independent half-line
boundary values are introduced. The proof does not assert convergence
of \(\partial_u r_m\), which would require control of a different
derivative. It proves exactly the fixed-source convergence of the
normal and tangential parts of the original first derivative, with
their sum and arithmetic source fixed at every degree.
